% STATUS: DRAFT
\chapter{Spectral theorem and unbounded operators}\label{ch:spectral}

\begin{physmotiv}
Every physicist writes $H=\sum_nE_n\ket n\!\bra n$ or $x=\int x\,\ket x\!\bra x\,\dd x$ and moves on. The spectral theorem is the statement that this is legitimate in general, with the sum and the integral unified into a single object, a \emph{projection-valued measure}. It also tells us what ``$f(H)$'' means for any function $f$, which we will use constantly (modular operators are $\Delta^{it}=\ee^{-\ii tK}$ and $\rho^{it}$). The second half of the chapter explains why $x$, $p$, $H$ are only defined on a \emph{domain}, and why ``Hermitian'' is not the same as ``self-adjoint''.
\end{physmotiv}

\section{The spectral theorem for bounded operators}

For a self-adjoint matrix $A$ with eigenprojections $P_\lambda$ we have $A=\sum_\lambda\lambda P_\lambda$. The general version replaces the sum with an integral over a projection-valued measure.

\begin{definition}
A \emph{projection-valued measure} (PVM) on $\R$ assigns to each Borel set $\Omega\subset\R$ a projection $E(\Omega)$ with $E(\R)=\id$, $E(\Omega_1\cap\Omega_2)=E(\Omega_1)E(\Omega_2)$, and $E(\bigcup_k\Omega_k)=\sum_kE(\Omega_k)$ (strongly) for disjoint $\Omega_k$.
\end{definition}

Physically, $E(\Omega)$ is the yes/no question ``is the value of the observable in $\Omega$?'', and $\inner\psi{E(\Omega)\psi}$ is the Born probability.

\begin{theorem}[spectral theorem]\label{thm:spectral}
Every self-adjoint operator $A$ (bounded or not) has a unique PVM $E_A$ supported on $\spec(A)$ with
\begin{equation}
A=\int_{\spec A}\lambda\,\dd E_A(\lambda).
\end{equation}
\end{theorem}

Examples: for $x$ on $L^2(\R)$, $E_x(\Omega)$ is multiplication by the indicator function $\chi_\Omega(x)$; for a matrix, $E(\Omega)=\sum_{\lambda\in\Omega}P_\lambda$. For the free-particle Hamiltonian the spectrum is $[0,\infty)$ and is purely continuous: there are no eigenvectors in $L^2$. ``Generalized eigenvectors'' $\ket p$ are not normalizable; the PVM is the rigorous replacement.

\subsection*{Functional calculus}
For any Borel function $f$ one defines
\begin{equation}
f(A)=\int f(\lambda)\,\dd E_A(\lambda).
\end{equation}
This is bounded iff $f$ is bounded on the spectrum, and $f\mapsto f(A)$ is a $^*$-homomorphism from bounded Borel functions to $\BH$. Important cases for us:
\begin{itemize}
\item $\ee^{-\ii tH}$, the time evolution (Stone, below);
\item $\rho^{it}$ and $\log\rho$ for a positive operator $\rho$ (the modular flow, modular Hamiltonian);
\item $|A|=\sqrt{A\ad A}$, and the spectral projection $\chi_{[0,\infty)}(A)$ (used in \cref{ch:clpw} to restrict the observer's energy $q\ge0$);
\item $\mathrm{sign}(A)$, $\theta(A)$, etc.
\end{itemize}

\begin{keyidea}
A self-adjoint operator is the same thing as a PVM on $\R$. Functions of an operator are defined through the PVM, and the \emph{projections} $\chi_\Omega(A)$ are bounded operators that live in any von Neumann algebra containing $A$ (``affiliated'' to it, see below).
\end{keyidea}

\section{Unbounded operators and domains}

An unbounded operator $A$ is a linear map defined on a \emph{dense} subspace $D(A)\subset\HH$, its \emph{domain}. Two operators with the same formula but different domains are different operators.

\begin{definition}
The adjoint $A\ad$ has domain $D(A\ad)=\{\eta:\xi\mapsto\inner\eta{A\xi}\text{ is bounded on }D(A)\}$ and is defined by $\inner{A\ad\eta}{\xi}=\inner\eta{A\xi}$. $A$ is \emph{symmetric} if $A\subset A\ad$ (i.e.\ $\inner\eta{A\xi}=\inner{A\eta}{\xi}$ on $D(A)$, the physicist's ``Hermitian''), and \emph{self-adjoint} if $A=A\ad$, including equality of domains. $A$ is \emph{essentially self-adjoint} if its closure is self-adjoint.
\end{definition}

Only self-adjoint operators have a spectral theorem, a one-parameter unitary group (Stone, below), and a unique time evolution. Symmetric but not self-adjoint operators can have none, or many.

\begin{example}[momentum on an interval]\label{ex:interval}
On $L^2[0,1]$ let $p=-\ii\,\dd/\dd x$.
\begin{itemize}
\item On $D_{\min}=\{\psi\text{ abs.\ cont.},\ \psi(0)=\psi(1)=0\}$ it is symmetric (boundary terms vanish) but has \emph{deficiency indices} $(1,1)$ and a $U(1)$ family of self-adjoint extensions: $\psi(1)=\ee^{\ii\alpha}\psi(0)$ (twisted periodic boundary conditions, with spectrum $2\pi n+\alpha$). Different $\alpha$ are physically different operators, different physics, with the same formula.
\item On the half line $[0,\infty)$ the operator $-\ii\dd/\dd x$ with $\psi(0)=0$ is symmetric but has deficiency indices $(0,1)$ or $(1,0)$ depending on sign: it has \emph{no} self-adjoint extension. There is no observable ``momentum on the half line''.
\end{itemize}
\end{example}

\begin{whatwrong}
The product of two symmetric operators is not symmetric, the sum of two self-adjoint operators may be defined only on the intersection of domains (which could be $\{0\}$), and the relation $[x,p]=\ii$ holds only on a common dense invariant domain. Formal manipulations such as $\ee^{A+B}=\ee^A\ee^B\ee^{-[A,B]/2}$ require care. The Weyl form (\cref{ch:weyl}) avoids these problems by using unitaries only.
\end{whatwrong}

A useful sufficient criterion (Nelson, Kato--Rellich): if $H_0$ is self-adjoint and $V$ is $H_0$-bounded with relative bound $<1$ then $H_0+V$ is self-adjoint on $D(H_0)$. Atomic Hamiltonians $-\nabla^2+V$ with Coulomb potential are essentially self-adjoint on $C_c^\infty$ (Kato). We will not need these in detail; the point is that the physical Hamiltonian is a \emph{choice} of self-adjoint operator.

\section{Stone's theorem}

\begin{theorem}[Stone]
There is a one-to-one correspondence between strongly continuous one-parameter unitary groups $U(t)$ and self-adjoint operators $H$, via $U(t)=\ee^{-\ii tH}=\int\ee^{-\ii t\lambda}\dd E_H(\lambda)$. $H$ is bounded iff $U$ is norm-continuous. $D(H)$ is the set of $\xi$ for which $t\mapsto U(t)\xi$ is differentiable.
\end{theorem}
Thus time evolution, translations, rotations, \dots\ determine self-adjoint generators, and a one-parameter group of unitaries is the same data as a self-adjoint generator. The modular flow $\Delta^{it}$ of \cref{ch:tomita} is the one-parameter group generated by $-\log\Delta$.

\section{Affiliated operators}

An unbounded operator cannot belong to a von Neumann algebra $\M\subset\BH$ (whose elements are bounded), but it can be \emph{affiliated} with it: $A$ is affiliated with $\M$ if $UAU\ad=A$ for every unitary $U\in\M'$. Equivalently, all spectral projections $\chi_\Omega(A)$ lie in $\M$. This is the right way to speak of ``the unbounded observable $x$'' inside a von Neumann algebra. Two examples to keep in mind: the observer's energy $q\ge0$ of \cref{ch:add_observer} will be an unbounded operator affiliated with the algebra of the observer; and the modular operator $\Delta$ of \cref{ch:tomita} is in general affiliated with neither $\M$ nor $\M'$, although conjugation by $\Delta^{it}$ maps $\M$ to itself.

\section*{Exercises}
\begin{exercise}
Show that the multiplication operator $x$ on $L^2(\R)$ with domain $\{\psi:\int x^2|\psi|^2<\infty\}$ is self-adjoint, and find $E_x(\Omega)$. What is the domain of $x^2$? Show $\psi(x)=(1+x^2)^{-1/2}$ lies in $L^2$ but not in $D(x)$.
\end{exercise}
\begin{exercise}
In \cref{ex:interval}, compute the deficiency spaces $\ker(p\ad\mp\ii)$ for the interval and the half line, and verify the claims about self-adjoint extensions. For the interval, find the spectrum of each extension.
\end{exercise}
\begin{exercise}
Let $\rho$ be a positive invertible matrix. Using the functional calculus show that $\rho^{it}=\ee^{\ii t\log\rho}$ is unitary, that $t\mapsto\rho^{it}$ is a one-parameter group and identify its generator.
\end{exercise}
\begin{exercise}
Show that if $A$ is self-adjoint and $A\ge0$ then $\ee^{-sA}$, $s>0$, is a bounded positive operator with norm $\le1$, although $A$ is unbounded. (This is the Euclidean heat kernel.) What happens for $\ee^{+sA}$?
\end{exercise}
